% [原文 p.93]
把 $H^0(X;\Z)$ 定义为由所有连续映射 $X\to\Z$ 组成的环。此时，作为群的直和分解
\[
K(X)=K_1(X)\oplus H^0(X;\Z)
\]
由同态 $\rank$ 得到。容易看出，自然同态
\[
H^0(X;\Z)\longrightarrow K_1(X)
\]
只有零映射。因此，若把 $\operatorname{Op}(H^0(\Z))$ 取为由 $\Z$ 到 $\Z$ 的所有映射组成的环，则 $\operatorname{Op}(K)=\operatorname{Op}(K,K)$ 与 $\operatorname{Op}(K_1,K)$ 的差别只在 $\operatorname{Op}(H^0(\Z))$。这样，定理 3.1.7 本质上给出了 $\operatorname{Op}(K)$ 的完整描述。

那么，转到对 $K(X)$ 的有限性条件的考察。先考虑 $H^0(X;\Z)$。

\begin{proposition*}{3.1.8}
下面两个条件等价：
\begin{enumerate}[label=(\Alph*),leftmargin=2.2em,itemsep=2pt]
\item $H^0(X;\Z)$ 是 Noether 环。
\item $H^0(X;\Z)$ 是有限生成 $\Z$ 模。
\end{enumerate}
\end{proposition*}

\begin{proof}
由 (B) 推 (A) 是显然的。于是设 $H^0(X;\Z)$ 为 Noether 环。假定能够找到 $X$ 的连通分支（既开又闭的集合）的严格递减无限链
\[
X=X_0\supsetneq X_1\supsetneq\cdots\supsetneq X_n\supsetneq X_{n+1}\supsetneq\cdots.
\]
此时，对每个 $n$，都能找到连续映射 $f_n\colon X\to\Z$，使得
\[
f_n(X_{n+1})=0,\qquad f_n(X_n-X_{n+1})=1
\]
% [原文 p.94]
成立。把 $H^0(X;\Z)$ 的理想 $I$ 取为：由对某个 $n$ 满足 $f(X_n)=0$ 的映射 $f\colon X\to\Z$ 组成。由于 $H^0(X;\Z)$ 是 Noether 环，$I$ 是有限生成的。因此存在 $N$，使得对一切 $f\in I$ 都有
\[
f(X_N)=0.
\]
然而，由于
\[
f_N\in I,\qquad f_N(X_N)\ne 0,
\]
这矛盾。因此，$X$ 只有有限个连通分支，从而由连通的 $X_i$ 表为
\[
X=\prod_{i=1}^{n}X_i.
\]
于是 $H^0(X;\Z)$ 同构于 $\Z^n$。
\end{proof}

% [Tutor Note: 这里的 $\\prod$ 在原书中表示各连通分支的不交并，而非拓扑积；$H^0$ 在每个连通分支上取 $\\Z$，故不交并给出直积 $\\Z^n$，两种读法在此恰好一致。]

那么，转到 $K(X)$。

\begin{proposition*}{3.1.9}
下面两个条件等价：
\begin{enumerate}[label=(\Alph*),leftmargin=2.2em,itemsep=2pt]
\item $K(X)$ 是 Noether 环。
\item $K(X)$ 是有限生成 $\Z$ 模。
\end{enumerate}
\end{proposition*}

\begin{proof}
这一次同样，若假定 (A)，则 $H^0(X;\Z)$ 是 $K(X)$ 的商环，从而是 Noether 环。于是由命题 3.1.8，$H^0(X;\Z)$ 是有限 $\Z$ 模。其次，$K_1(X)$ 是由 $K(X)$ 的幂零元组成的理想（系 3.1.6）。而 $K(X)$ 是 Noether 环，所以 $K_1(X)$ 是幂零理想。为便于观察，记 $I=K_1(X)$。此时，对某个 $n$ 有 $I^n=0$，并且 $I^m/I^{m+1}$（$m=0,1,\dots,n-1$）都是 $K/I=H^0(X;\Z)$ 上的有限生成模。因此，$K(X)$ 是有限生成 $H^0(X;\Z)$ 模，从而是有限生成 $\Z$ 模。
\end{proof}

$K(X)$ 为有限生成 $\Z$ 加群的空间 $X$ 的例子有胞腔复形。

% [原文 p.95]
这里，用由满足 $\sum_{j=1}^{k}i_j\ge n$ 且 $x_i\in K_1(X)$ 的单项式
\[
\gamma^{i_1}(x_1)\gamma^{i_2}(x_2)\cdots\gamma^{i_k}(x_k)
\]
全体生成的子群 $K_n^\gamma(X)$，来定义 $K(X)$ 的滤过。由于 $\gamma^1(x)=x$，有 $K_1^\gamma=K_1$。当 $x\in K_n^\gamma(X)$ 时，称 $x$ 的 $\gamma$ 滤过不小于 $n$，写作 $F_\gamma(x)\ge n$。

\begin{proposition*}{3.1.10}
设 $K(X)$ 为有限生成 $\Z$ 模。此时，对某个 $n$ 有
\[
K_n^\gamma(X)=0
\]
成立。
\end{proposition*}

\begin{proof}
设 $x_1,\dots,x_n$ 为 $K_1(X)$ 的生成元，$N_j=N(x_j)$ 为命题 3.1.5 给出的整数。由式
\[
\gamma_t(a+b)=\gamma_t(a)\gamma_t(b),
\]
只需证明存在 $N$，使得 $\gamma^i(x_j)$ 的单项式只要重量之和大于 $N$ 就等于零。然而，若取 $N=\sum_{j=1}^{n}N_j$，则这样的任意单项式对某个 $j$ 必有 $\gamma^i(x_j)$ 的重量大于 $N_j$。于是由命题 3.1.5，这个单项式为零。
\end{proof}

\begin{corollary*}{3.1.11}
假定 $K(X)$ 为有限生成 $\Z$ 模。此时 $\dim_\gamma X$ 有限。
\end{corollary*}

再列举算子 $\gamma^i$ 的一些进一步的性质。

\begin{proposition*}{3.1.12}
若 $V$ 是 $n$ 维向量丛，则
\[
\lambda_{-1}[V]=(-1)^n\gamma^n([V]-n).
\]
因此，$\tilde K^*(X^V)$ 是自由 $K^*(X)$ 模，并由 $\gamma^n([V]-n)$ 生成。
\end{proposition*}

% [原文 p.96]
\begin{proposition*}{3.1.13}
存在多项式 $P_i$、$Q_{ij}$，使得对一切 $x$、$y$ 有
\begin{align*}
\gamma^i(xy)&=P_i(\gamma^1(x),\gamma^1(y),\gamma^2(x),\gamma^2(y),\dots,\gamma^i(x),\gamma^i(y)),\\
\gamma^i(\gamma^j(x))&=Q_{ij}(\gamma^1(x),\dots,\gamma^{i+j}(x)).
\end{align*}
\end{proposition*}

这些证明留给读者。用分裂原理应当可以简单验证。

\section{Adams 算子}
\takeaway{Adams 算子 $\psi^k$ 由 $\psi_t(x)=\psi^0(x)-t\frac{d}{dt}\log\lambda_{-t}(x)$ 定义，它保持加法与乘法并把 K 群按特征值 $k^i$ 分裂。由此可重建 Betti 数与欧拉示性数，并得到 Hopf 不变量的整性定理 3.2.3 与 3.2.10。}

这里，我们想关注具有特别良好性质的算子。它由 J. F. Adams 引入。定义 $\psi^0(x)=\rank x$。在环 $K(X)[[t]]$ 中，用
\[
\psi_t(x)=\psi^0(x)-t\frac{d}{dt}\bigl(\log\lambda_{-t}(x)\bigr)
\]
定义 $\psi_t(x)=\sum_{i=0}^{\infty}t^i\psi^i(x)$。注意，这个幂级数的系数全是整数，所以这个定义有意义。

% [Tutor Note: 这里用的是 $\\lambda_{-t}$ 而非 $\\lambda_t$，并且前面带 $-t\\frac{d}{dt}$：这两处符号配合起来，才使线丛情形给出 $\\psi^k(x)=x^k$（命题 3.2.1 第 2 项）；写成 $\\lambda_t$ 会把符号弄反。]

\begin{proposition*}{3.2.1}
对任意 $x,y\in K(X)$，下面各项成立。
\begin{enumerate}[label=\arabic*),leftmargin=2.2em,itemsep=2pt]
\item 对一切 $k$，$\psi^k(x+y)=\psi^k(x)+\psi^k(y)$。
\item 若 $X$ 为线丛，则 $\psi^k(x)=x^k$。
\item 由前两项，算子 $\psi^k$ 被唯一确定。
\end{enumerate}
\end{proposition*}

% [Tutor Note: 原日译本 p.96 此处（以及本命题第 2 项、证明第二段）印作大写的 $X$：$X$ 是空间记号，此处应指元素 $x$ 为线丛的类，大写 $X$ 疑为笔误；正文照印。]

\begin{proof}
由于 $\psi_t(x+y)=\psi_t(x)+\psi_t(y)$，对每个 $k$ 有 $\psi^k(x+y)=\psi^k(x)+\psi^k(y)$。

若 $X$ 为线丛，则 $\lambda_{-t}(x)=1-tx$，因此
\begin{align*}
\frac{d}{dt}\bigl(\log(1-tx)\bigr)&=\frac{-x}{1-tx}\\
&=-x-tx^2-t^2x^3-\cdots
\end{align*}
成立。从而 $\psi_t(x)=1+tx+t^2x^2+\cdots$ 成立。

% [原文 p.97]
主张 3) 由分裂原理推出。
\end{proof}

\begin{proposition*}{3.2.2}
对任意 $x,y\in K(X)$，下面各项成立。
\begin{enumerate}[label=\arabic*),leftmargin=2.2em,itemsep=2pt]
\item 对一切 $k$，$\psi^k(xy)=\psi^k(x)\psi^k(y)$。
\item 对一切 $k,\ell$，$\psi^k(\psi^\ell(x))=\psi^{k\ell}(x)$。
\item 若 $p$ 为素数，则 $\psi^p(x)\equiv x^p \bmod p$。
\item 若 $u\in\tilde K(S^{2n})$，则对一切 $k$ 有 $\psi^k(u)=k^nu$。
\end{enumerate}
\end{proposition*}

\begin{proof}
主张 1) 与 2) 由命题 3.2.1 和分裂原理立刻推出。仍由分裂原理，设 $f$ 为整系数多项式，则有 $\psi^p(x)=x^p+pf(\lambda^1(x),\dots,\lambda^p(x))$。最后，若 $h$ 是 $\tilde K(S^2)$ 的生成元，则 $\psi^k(h)=kh$。而 $S^{2n}=S^2\wedge\cdots\wedge S^2$，且 $\tilde K(S^{2n})$ 由 $h\otimes\cdots\otimes h$ 生成，所以主张 4) 由主张 1) 推出。
\end{proof}

下面给出 Adams 算子 $\psi^k$ 的应用。对任意映射 $f\colon S^{4n-1}\to S^{2n}$，如下定义 Hopf 不变量 $H(f)$。设 $X_f$ 为 $f$ 的映射锥。设 $i\colon S^{2n}\to X_f$ 为包含映射，$j\colon X_f\to S^{4n}$ 为把 $S^{2n}$ 收缩为一点的映射。设 $u$ 为 $\tilde K(S^{4n})$ 的生成元。由对 $(X_f,S^{2n})$ 的正合列，存在元素 $x\in\tilde K(X_f)$，使 $i^*(x)$ 生成 $\tilde K(S^{2n})$。$\tilde K(X_f)$ 是由 $x$ 与 $y=j^*(u)$ 生成的自由 Abel 群。由于 $(i^*(x))^2=0$，对某个 $H$ 有 $x^2=Hy$。把这个整数 $H$ 定义为 $f$ 的 Hopf 不变量。显然，除正负号外，$H(f)$ 是良定义的。下面的定理由 J. F. Adams 用上同调方法首先证明。

% [Tutor Note: 为何 $x^2$ 必是 $y$ 的整数倍：$\\tilde K(X_f)$ 以 $x,y$ 为自由 Abel 基，而 $i^*(x^2)=(i^*x)^2=0$ 说明 $x^2$ 落在 $\\tilde K(X_f,S^{2n})$ 中，这个子群由 $y=j^*(u)$ 生成，故 $x^2=Hy$。Hopf 不变量的整性是这一构造的全部要点。]

\begin{theorem*}{3.2.3}
若 $H(f)$ 为奇数，则 $n$ 是 $1,2,4$ 之一。
\end{theorem*}

\begin{proof}
设 $\psi^2(x)=2^nx+ay$，$\psi^3(x)=3^nx+by$。由于 $\psi^2(x)\equiv x^2\bmod 2$，$a$ 是奇数。又 $\psi^k(y)=j^*(\psi^k(u))=k^{2n}y$ 成立。因此，
\begin{align*}
\psi^6(x)&=\psi^3(\psi^2(x))=6^nx+(2^nb+3^{2n}a)y,\\
\psi^6(x)&=\psi^2(\psi^3(x))=6^nx+(2^{2n}b+3^na)y
\end{align*}
成立。于是 $2^nb+3^{2n}a=2^{2n}b+3^na$，即 $2^n(2^n-1)b=3^n(3^n-1)a$。

% [原文 p.98]
由于 $a$ 是奇数，$2^n$ 整除 $3^n-1$。因此由初等数论，$n$ 是 $1,2,4$ 之一。
\end{proof}

当 $n$ 为 $1,2,4$ 时，把 $S^{4n-1}$ 看作复数、四元数或 Cayley 数（八元数）上二维向量空间中非零向量组成的子空间，把 $S^{2n}$ 看作复射影直线、四元数射影直线或 Cayley 数射影直线，便确定了 Hopf 映射。这些 Hopf 不变量都是 1。验证这一点留给读者。

\begin{proposition*}{3.2.4}
设 $x\in K(X)$ 满足 $F_\gamma(x)\ge n$。此时，对任意 $k$，
\[
F_\gamma(\psi^k(x)-k^nx)\ge n+1
\]
成立。
\end{proposition*}

\begin{proof}
若 $n=0$，则
\[
\psi^k(x)=\psi^k(\rank x+x_1)=\rank x+\psi^k(x_1)
\]
成立。此时 $x_1$ 与 $\psi^k(x_1)$ 都含于 $K_1(X)$。因此
\[
\psi^k(x)-x=\psi^k(x_1)-x_1\in K_1(X)=K_1^\gamma(X)
\]
成立。其次设 $n>0$。由于 $\psi^k$ 是环同态，只需证明：以 $\psi^k\in\operatorname{Op}(K)$、$\gamma^n\in\operatorname{Op}(K_1,K)$ 记时，合成 $\psi^k\circ\gamma^n-k^n\gamma^n$ 等于 $\gamma^i$ 的一个多项式，且各项的重量不小于 $n+1$。与定理 3.1.7 同样，有同构
\[
\Z[[\gamma^1,\dots]]\cong\operatorname{Op}(K_1,K)\cong\lim_{m}\Z[x_1,\dots,x_m]^{S_m}.
\]
此时，$\gamma^i$ 对应 $x_j$ 的 $i$ 次初等对称多项式 $\sigma_i$。现在，
\[
\psi^k(x_i)=(1+x_i)^k-1
\]
成立，所以
% [原文 p.99]
\begin{align*}
\psi^k(\sigma_n(x_1,\dots))&=\sigma_n((1+x_1)^k-1,\dots)\\
&=k^n\sigma_n(x)+f
\end{align*}
成立。其中 $f$ 是 $\sigma_i$ 的多项式，重量不小于 $n+1$。由前面所述的同构，$\psi^k\circ\gamma^n$ 对应 $\psi^k(\sigma_n)$，故命题得证。
\end{proof}

反复使用命题 3.2.4，便得到下面的推论。

\begin{corollary*}{3.2.5}
若 $K_{n+1}^\gamma(X)=0$，则对非负整数的任意数列 $k_0,k_1,\dots,k_n$，
\[
\prod_{m=0}^{n}\bigl(\psi^{k_m}-(k_m)^m\bigr)=0
\]
成立。
\end{corollary*}

% [Tutor Note: 这个乘积公式的来历：命题 3.2.4 说每个因子 $\\psi^k-k^n$ 把 $\\gamma$ 滤过至少提高 1；从 $F_<本地路径> 出发乘上 $m=0,1,\\dots,n$ 这 $n+1$ 个因子后落在 $K_{n+1}^\\gamma(X)=0$ 中，故整体为零算子。]

由命题 3.1.10，特别地，当 $K(X)$ 为有限生成 $\Z$ 模时，总能应用系 3.2.5。

注意，$\psi^k$ 作为线性变换作用于向量空间 $K(X)\otimes\Q$。在系 3.2.5 中，若对一切 $m$ 取 $k_m=k$，则可知在 $K(X)\otimes\Q$ 上
\[
\prod_{m=0}^{n}(\psi^k-k^m)=0
\]
成立。因此，各 $\psi^k$ 的特征值是不超过 $k^n$ 的 $k$ 的幂。用 $V_{k,i}$ 表示特征值 $k^i$ 所对应的 $\psi^k$ 的特征空间（也可能 $V_{k,i}=0$）。此时，若 $k>1$，则 $K(X)\otimes\Q$ 的恒等算子 1 有正交分解
\[
1=\sum_i\Pi_i,\qquad
\Pi_i=\prod_{m\ne i}(\psi^k-k^m)/(k^i-k^m).
\]
因此，$K(X)\otimes\Q$ 是 $V_{k,i}$ 的直和。这里，在系 3.2.5 中，

% [原文 p.100]
取 $k_i=\ell$、$k_m=k$（$i\ne m$），便知
\[
(\psi^\ell-\ell^i)V_{k,i}=0,
\]
所以 $V_{k,i}\subset V_{\ell,i}$。因此，下面的命题得证。

\begin{proposition*}{3.2.6}
假定 $K(X)$ 具有有限的 $\gamma$ 滤过。用 $V_{k,i}$ 表示 $K(X)\otimes\Q$ 上特征值 $k^i$ 所对应的 $\psi^k$ 的特征空间。此时，若 $k,\ell>1$，则
\[
V_{k,i}=V_{\ell,i}
\]
成立。
\end{proposition*}

由命题 3.2.6，当 $k>1$ 时子空间 $V_{k,i}$ 不依赖于 $k$，所以可以用与 $k$ 无关的记号表示。于是把 $V_{k,i}$ 写作 $H^{2i}(X;\Q)$，称之为空间 $X$ 的 $2i$ 次 Betti 群（Betti group）\footnote{译注：一般多称为上同调群。}。由命题 3.2.4 可知，特征值 $k^0=1$ 只在 $H^0(X;\Z)\otimes\Q$ 中出现。因此，在
\[
H^0(X;\Z)\otimes\Q=H^0(X;\Q)
\]
的意义下，这个记号是一致的。定义奇数次的 Betti 群为
\[
H^{2m+1}(X;\Q)=H^{2m+2}(SX^+;\Q).
\]
其中 $X^+=X\cup\{*\}$，$S$ 表示约化悬挂。若这些空间的上同调有限维，则令
\[
B_k=\dim_\Q H^k(X;\Q),
\]
并把欧拉示性数 $E(X)$ 定义为
\[
E(X)=\sum(-1)^kB_k=\dim_\Q(K^0(X)\otimes\Q)-\dim_\Q(K^1(X)\otimes\Q).
\]

% [Tutor Note: 这里把 $\\psi^k$ 的特征空间直接命名为 $H^{2i}(X;\\Q)$，只给出偶数次的上同调；奇数次靠 $\\psi^k$ 看不到，才用悬挂 $H^{2m+1}(X;\\Q)=H^{2m+2}(SX^+;\\Q)$ 人为定义。故这套「Betti 群」与通常奇异上同调的一致性并非平凡事实。]

% [原文 p.101]
有了这个定义，当 Künneth 公式可以适用时，可以导出
\[
E(X\times Y)=E(X)E(Y).
\]
注意这一点。

下面的命题，只是用这里引入的记号把命题 3.2.4 改写了一遍。

\begin{proposition*}{3.2.7}
\[
K_n^\gamma(X)\otimes\Q=\sum_{m\ge n}H^{2m}(X;\Q),
\]
因此
\[
\{K_n^\gamma(X)/K_{n+1}^\gamma(X)\}\otimes\Q\cong H^{2n}(X;\Q)
\]
成立。
\end{proposition*}

对 $\tilde K(S^2)$ 的生成元 $u$ 有 $\psi^k(u)=ku$，所以
\[
\psi^k\beta(x)=k\beta\psi^k(x)
\]
成立。其中 $\beta\colon K(X)\to K^{-2}(X)$ 是周期性定理的同构映射。因此，$\beta$ 诱导同构
\[
H^{2m}(X;\Q)\cong H^{2m+2}(S^2X^+;\Q).
\]
由定义奇数次 Betti 群的方法可知，对一切 $k$，
\begin{equation}
H^k(X;\Q)\cong H^{k+1}(SX^+;\Q) \tag{3.2.8}
\end{equation}
成立。

% [Tutor Note: 原书中公式编号与陈述编号共用「节.序号」计数器：命题 3.2.7 之后直接是公式 (3.2.8)，再接命题 3.2.9。抄录时不要误补出一个「命题 3.2.8」。]

这里，对 $(X,A)$ 的 K 群的正合列取与 $\Q$ 的张量积，在 $\psi^k$ 下分解，并使用式 (3.2.8)，便得到下面的命题。

\begin{proposition*}{3.2.9}
设 $A\subset X$。若 $K^*(X)$、$K^*(A)$ 都是有限 $\Z$ 模，则正合列
% [原文 p.102]
\[
\cdots\to K^{i-1}(A)\xrightarrow{\ \delta\ }K^i(X,A)\to K^i(X)\to K^i(A)\xrightarrow{\ \delta\ }\cdots
\]
导出正合列
\[
\cdots\to H^{i-1}(A;\Q)\xrightarrow{\ \delta\ }H^i(X,A;\Q)\to H^i(X;\Q)\to H^i(A;\Q)\xrightarrow{\ \delta\ }\cdots.
\]
\end{proposition*}

下面给出算子 $\psi^k$ 的另一个应用。由于 $P_n(\C)/P_{n-1}(\C)$ 是球面 $S^{2n}$，对每个 $k$ 都有从 $S^{2n}$ 到 $P_{n+k}(\C)/P_{n-1}(\C)$ 的包含映射。我们想知道，$n$ 与 $k$ 取什么值时 $S^{2n}$ 成为 $P_{n+k}(\C)/P_{n-1}(\C)$ 的收缩。也就是说，想知道在 $S^{2n}$ 上成为恒等映射的映射 $f\colon P_{n+k}(\C)/P_{n-1}(\C)\to S^{2n}$ 何时存在。下面求这样的 $f$ 存在时 $n$ 与 $k$ 的必要条件。

\begin{theorem*}{3.2.10}
设有收缩
\[
f\colon P_{n+k}(\C)/P_{n-1}(\C)\longrightarrow P_n(\C)/P_{n-1}(\C)=S^{2n}.
\]
此时，$\left(\dfrac{\log(1+x)}{x}\right)^n$ 中 $x^i$（$i\le k$）的系数全是整数。
\end{theorem*}

\begin{proof}
设 $\xi$ 为 $P_{n+k}(\C)$ 上的典范线丛，$x=\xi-1$。于是 $K(P_{n+k}(\C))$ 是以 $x^s$（$0\le s\le n+k$）为生成元的自由 Abel 群，而 $K(P_{n+k}(\C),P_{n-1}(\C))$ 可以等同于由 $x^s$（$n\le s\le n+k$）生成的部分群。在 $K(P_{n+k}(\C))\otimes\Q$ 中，令 $y=\log(1+x)$，则 $\xi=e^y$。此时，
\[
e^{ry}=\xi^r=\psi^r(e^y)=e^{\psi^r(y)}
\]
成立，所以 $\psi^r(y)=ry$。因此，对 $n\le s\le n+k$，$H^{2s}(P_{n+k}(\C)/P_{n-1}(\C);\Q)$ 是由 $y^s$ 生成的 1 维空间。这里，设 $u$ 生成 $\tilde K(S^{2n})$，
\end{proof}

% 块尾环境：定理 3.2.10 的证明在 p.103 由 ch03c 以「证明（续）」接续
% CHUNK ch03b DONE
